\section{A First Example: Array Initialization}\label{sec:init}

This section follows one proof from the program to the closed goal, step by step, as the engine performs it. The
program fills the first $n$ cells of an array with a value; a Hoare-style proof needs the invariant
$\forall k.\ 0 \le k < i \Rightarrow a[k] = v$. Here, nothing is written for the program: the proof uses the lemmas
of the database only.
\begin{lstlisting}
method init(n: int, v: int) returns (a: array<int, int>)
  requires 0 <= n
  ensures forall k: int :: 0 <= k < n ==> a[k] == v
{
  var i: int;
  i := 0;
  while i < n { a[i] := v; i := i + 1; }
}
\end{lstlisting}
The run is \texttt{./Cuvee.sh -teq -prove examples/dafny/teq/init.dfy}; its output has no goal, i.e.\ the goal is
proved.

\subsection{The verification condition}
The loop modifies $a$ and $i$. The index reads nothing, the array reads the index, so the two are separate
components (Section~\ref{sec:teq}), the array after the index. With \texttt{-print dfy} in place of
\texttt{-prove}, and the fresh variables numbered by the tool, the goal is
\begin{align*}
&\forall n, v, a.\ 0 \le n \Rightarrow \forall a', i', n_4{:}\mathsf{Nat}, K, I, A.\\
&\quad K = \cnt(0, 1, n_4)\ \wedge\ I = \tr{0}{\bl{K}}{i, k \mapsto i + 1}\\
&\quad \wedge\ A = \tr{a}{\bl{K}, \bl{I}}{a', k, i \mapsto a'[i := v]}\
 \wedge\ a' = \cur(A)\ \wedge\ i' = \cur(I)\\
&\quad \wedge\ \al{\bl{K}, \bl{A}, \bl{I}}{k, a', i \mapsto i < n}\ \wedge\ 0 \le \intf(n_4)\\
&\quad \Rightarrow \neg(i' < n) \Rightarrow \forall k.\ 0 \le k < n \Rightarrow a'[k] = v.
\end{align*}
Every construct of the approach appears once. The number $n_4$ of iterations is a fresh natural, about which nothing
is known but what the iteration constraint and the exit say. $K$ is the counter of the iterations. The trace $I$ of
the index starts at $0$ and steps by one; the trace $A$ of the array has the history $\bl{I}$ as a further
dependency, and in iteration $k$ its update reads the index of that iteration through the bound variable $i$, i.e.\
$i = \ear(\natf(k), \bl{I})$. The exit values are the current values of the traces, and the iteration constraint
states the loop test for every earlier iteration. Apart from the updates and the test, the condition is the schema
that every loop gets; it contains no invariant.

\subsection{The lemmas}
The proof uses these lemmas of the database (\texttt{resources/cuvee/teq.dfy}):
\begin{lstlisting}
lemma {:rw} counter_def(i: int, c: int, n: Nat)
  ensures trace(i, for ⌊counter(0, 1, n)⌋ do x, k :: x + c) == counter(i, c, n)
lemma {:rw} current_counter(i: int, c: int, n: Nat)
  ensures current(counter(i, c, n)) == i + integer(n) * c
lemma {:rw} earlier_counter(k: Nat, i: int, c: int, n: Nat)
  ensures earlier(k, ⌊counter(i, c, n)⌋) == i + integer(k) * c
lemma {:rw} invs_counter_upper_nonnegative(m: int, n: Nat)
  requires 0 <= m
  ensures invs(for ⌊counter(0, 1, n)⌋ check k :: k < m) <==> integer(n) <= m

function {:fold} stores<a, b>(a: array<a, b>, keys: List<a>, values: List<b>): Trace<array<a, b>>
{ trace(a, for keys, values do b, k, v :: b[k := v]) }
lemma {:rw} maps_const<a>(x: a, i: int, c: int, n: Nat)
  ensures maps(for ⌊counter(i, c, n)⌋ yield v :: x) == ⌊consts(n, x)⌋
lemma {:rw} current_stores<a, b>(k: a, a: array<a, b>, keys: List<a>, values: List<b>)
  ensures current(stores(a, keys, values))[k]
    == if occurs(k, keys) then earlier(indexof(k, keys), values) else a[k]

lemma {:rw} occurs_counter_up(j: int, i: int, n: Nat)
  ensures occurs(j, ⌊counter(i, 1, n)⌋) <==> i <= j < i + integer(n)
lemma {:rw} indexof_counter_up(j: int, i: int, n: Nat)
  requires i <= j < i + integer(n)
  ensures indexof(j, ⌊counter(i, 1, n)⌋) == natural(j - i)
lemma {:rw} earlier_consts<a>(k: Nat, n: Nat, x: a)
  ensures earlier(k, ⌊consts(n, x)⌋) == x
\end{lstlisting}
They are of three kinds. \emph{Introduction} lemmas recognize a shape and name it: a trace that adds a constant
is a counter, and a trace that stores values at keys is $\stores$, a function marked \texttt{\{:fold\}}, whose body
is rewritten \emph{to} the function. \emph{Acceleration} lemmas give the closed form of a named shape: the last
value of a counter, its $k$-th value, the value of a store sequence at a key, which is the last value stored there
or the original one. \emph{Bookkeeping} lemmas relate the functions on lists: $\occurs$, $\indexof$, the constant list
$\consts$, and $\mapc{\ldots}{\ldots}$, which the engine introduces itself when a lemma needs a list that the
program computes only implicitly. All of them are proved in Lean (Section~\ref{sec:engine:db}).

\subsection{The proof}
The engine simplifies the goal under a context of facts, innermost terms first, and applies the first rule whose
left-hand side matches and whose side conditions hold (Section~\ref{sec:engine}). The steps below are those of
\texttt{-debug rewrite}, cleaned up.

\paragraph{The index is a counter.} The trace of the index is matched by \lemma{counter\_def} as it stands,
$\tr{0}{\bl{K}}{i, k \mapsto i + 1} \rw \cnt(0, 1, n_4)$. From here on, $I = \cnt(0, 1, n_4)$ is a fact, and a
lemma that asks for a counter where the goal has $\bl{I}$ matches through this equation.

\paragraph{The index read by the array is the iteration number.} To simplify the update of $A$, the engine builds the
context of the combinator: $k$ is a value of $\bl{K}$, so $0 \le k < \intf(n_4)$, and $i$ is the value of $\bl{I}$ at
position $k$. That read is simplified in the same context,
$\ear(\natf(k), \bl{I}) \rw 0 + \intf(\natf(k)) \cdot 1 \rw k$, so in every iteration the index is $k$, and the update
of $A$ becomes $a'[k := v]$. ``The loop stores at $i$'' has become ``the loop stores at $0, 1, \ldots$'', and the
unused dependency $\bl{I}$ is dropped.

\paragraph{The array is a sequence of stores.} The trace $\tr{a}{\bl{K}}{a', k \mapsto a'[k := v]}$ is matched by the
body of $\stores$ with $\bl{K}$ as the keys and, as the values, a list that the program never names: the engine
synthesizes it as the map of the stored value over the iterations,
\begin{align*}
  \tr{a}{\bl{K}}{a', k \mapsto a'[k := v]} &\rw \stores(a, \bl{K}, \mapc{\bl{K}}{k \mapsto v}),\\
  \mapc{\bl{K}}{k \mapsto v} &\rw \bl{\consts(n_4, v)}.
\end{align*}
The final array is $\cur(\stores(a, \bl{K}, \bl{\consts(n_4, v)}))$, and the final index
$\cur(\cnt(0, 1, n_4)) \rw \intf(n_4)$.

\paragraph{The number of iterations is $n$.} The iteration constraint has become $\al{\bl{K}}{k \mapsto k < n}$, which
\lemma{invs\_counter\_upper\_nonnegative} turns into $\intf(n_4) \le n$, its side condition $0 \le n$ being the
precondition. The exit gives $n \le \intf(n_4)$, and the arithmetic of the context concludes $\intf(n_4) = n$.

\paragraph{The postcondition.} $a'$ is a variable whose definition is $\cur(A)$, with $A$ the store sequence, and
\lemma{current\_stores} fires through both:
\begin{align*}
a'[k] &\rw \ite{\occurs(k, \bl{K})}{\ear(\indexof(k, \bl{K}), \bl{\consts(n_4, v)})}{a[k]}\\
\occurs(k, \bl{K}) &\rw 0 \le k < \intf(n_4), \text{ true for } 0 \le k < n = \intf(n_4)\\
\ear(\natf(k), \bl{\consts(n_4, v)}) &\rw v
\end{align*}
by \lemma{occurs\_counter\_up}, \lemma{indexof\_counter\_up} and \lemma{earlier\_consts}. The conclusion $a'[k] = v$ is
true, and the goal is dropped. Had the program stored $v$ at $2i$ or counted down, other lemmas of the same database
would have fired; the proof used nothing about the loop but what the database knows about counters, stores and
constant lists.
